%============================================================
% Chapter 17: Hall-Type Conditions and Greedy Feasibility
% STATUS: PLANNED  (generated by tools/outline.py; replace this file when the chapter is written)
%============================================================
\chapter{Hall-Type Conditions and Greedy Feasibility}
\label{ch:hall}

\begin{tcolorbox}[enhanced, colback=white, colframe=secondaryblue, boxrule=1pt, arc=2mm]
\textbf{\large\textcolor{primaryblue}{A greedy assignment exists exactly when no group of requests is too greedy itself: Hall's condition on prefixes, intervals and supplies.}}
\end{tcolorbox}

\chapterfacts{2000}{2600}{3}{6}{30}

\begin{plannedbox}[Planned Contents --- this chapter is not written yet]
\textbf{Part A --- Lessons}
\begin{enumerate}[leftmargin=*, label=\textbf{A\arabic*.}, itemsep=1pt]
    \item Hall's theorem and its interval and prefix forms
    \item Checking supply--demand feasibility with sorted prefixes
    \item Greedy assignment that succeeds iff Hall holds
    \item Deficiency: how many requests can be served
    \item Flows as the general tool (the graph book)
\end{enumerate}
\textbf{Part B --- Patterns} (with candidate problems, checked when the chapter is written)
\begin{enumerate}[leftmargin=*, label=\textbf{P\arabic*.}, itemsep=1pt]
    \item \textbf{Prefix Hall conditions}
    \item \textbf{Interval supplies}
    \item \textbf{Deficiency}
\end{enumerate}
\end{plannedbox}
